\chapter{Blockchain e Contratos Inteligentes}
\label{chap:blockchain-contratos}

\section{Objetivos de aprendizagem}
Ao final deste capítulo, o aluno deve ser capaz de:
\begin{itemize}[leftmargin=*]
  \item explicar o que caracteriza uma blockchain e por que ela difere de um banco de dados convencional;
  \item distinguir bloco, transação, estado, nó, consenso e contrato inteligente;
  \item compreender o papel da \textit{Ethereum Virtual Machine} (EVM), do custo computacional e do armazenamento on-chain;
  \item interpretar o papel de eventos, funções de leitura e funções de escrita em contratos;
  \item justificar quando vale a pena usar blockchain em uma arquitetura de identidade e controle de acesso.
\end{itemize}
\section{Do livro-razão distribuído ao contrato inteligente}
Uma blockchain pode ser entendida como um livro-razão distribuído no qual múltiplos nós mantêm uma visão compartilhada e verificável de transações e estado. Em vez de confiar em uma única autoridade central para armazenar e validar registros, a rede aplica regras de consenso e validação criptográfica sobre cada alteração \citep{nakamoto_bitcoin,ethereum_whitepaper}.

Embora o tema tenha ganhado popularidade com criptomoedas, o conceito relevante para este curso é mais amplo: blockchain é uma infraestrutura de registro verificável, com forte rastreabilidade e regras de atualização definidas por protocolo.

\section{Elementos básicos de uma blockchain}
Os componentes conceituais mais importantes são:
\begin{itemize}[leftmargin=*]
  \item \textbf{bloco}: agrupamento ordenado de transações e metadados;
  \item \textbf{transação}: operação solicitada por um participante da rede;
  \item \textbf{nó}: instância de software que valida, propaga ou armazena o estado da rede;
  \item \textbf{estado}: fotografia lógica atual dos dados relevantes para execução do protocolo;
  \item \textbf{consenso}: mecanismo pelo qual a rede concorda sobre a ordem e validade das alterações;
  \item \textbf{hash do bloco}: referência criptográfica que vincula blocos sucessivos.
\end{itemize}

Esse encadeamento por hash dificulta alteração retroativa invisível, porque uma mudança em um bloco compromete os vínculos criptográficos dos blocos seguintes.

\section{Por que blockchain não é apenas banco de dados}
Um banco de dados convencional costuma oferecer:
\begin{itemize}[leftmargin=*]
  \item alta flexibilidade de consulta;
  \item controle administrativo centralizado;
  \item atualização direta por operadores autorizados.
\end{itemize}

Uma blockchain, por sua vez, prioriza:
\begin{itemize}[leftmargin=*]
  \item verificabilidade compartilhada;
  \item histórico rastreável de transações;
  \item execução de regras padronizadas em ambiente distribuído;
  \item resistência a adulteração não consensual.
\end{itemize}

Essa diferença é essencial. Blockchain não substitui qualquer banco de dados; ela faz sentido quando auditoria, confiança distribuída, política verificável e integridade histórica são requisitos relevantes.

\section{Transações, contas e estado}
Em plataformas como Ethereum, a rede opera sobre um modelo de estado global. Cada bloco confirma transações que alteram esse estado. Em termos didáticos, uma transação pode:
\begin{itemize}[leftmargin=*]
  \item transferir valor;
  \item criar um contrato;
  \item invocar uma função de contrato;
  \item registrar um evento como efeito observável da execução.
\end{itemize}

As contas geralmente se dividem em:
\begin{itemize}[leftmargin=*]
  \item \textbf{contas externas}: controladas por chaves privadas de usuários ou serviços;
  \item \textbf{contas de contrato}: controladas pelo código do contrato implantado na rede.
\end{itemize}

\section{Consenso em alto nível}
O consenso é o mecanismo que permite à rede decidir qual é a próxima versão válida do livro-razão distribuído. O aluno não precisa dominar, neste momento, cada detalhe de implementação de prova de trabalho ou prova de participação. O ponto essencial é entender que:
\begin{itemize}[leftmargin=*]
  \item diferentes redes usam diferentes mecanismos de consenso;
  \item o consenso afeta desempenho, custo, segurança e finalização;
  \item a aplicação deve respeitar o tempo de confirmação e o modelo de confiança da rede usada.
\end{itemize}

Em ambientes didáticos, também é comum usar redes locais de desenvolvimento, nas quais a semântica do contrato é preservada, mas o consenso é simplificado para facilitar testes.

\section{Ethereum e a Ethereum Virtual Machine}
Ethereum se tornou uma das plataformas mais importantes para contratos inteligentes ao introduzir um ambiente geral de execução programável. Seu modelo permite armazenar estado e executar lógica determinística em todos os nós validadores relevantes da rede \citep{ethereum_whitepaper,wood_ethereum}.

A \textit{Ethereum Virtual Machine} (EVM) é a máquina abstrata que executa o bytecode dos contratos. Conceitualmente, ela:
\begin{itemize}[leftmargin=*]
  \item recebe instruções do contrato compilado;
  \item processa operações aritméticas, lógicas e de armazenamento;
  \item mantém um contexto de execução padronizado;
  \item garante que a execução seja determinística para todos os participantes.
\end{itemize}

Esse determinismo é crítico. Se dois nós executarem a mesma transação sobre o mesmo estado e chegarem a resultados diferentes, a rede perde consistência.

\section{Gas e custo computacional}
Em Ethereum e plataformas similares, a execução não é tratada como custo zero. Cada operação consome \textit{gas}, uma unidade abstrata associada ao esforço computacional e ao uso de recursos da rede \citep{wood_ethereum,evm_codes}.

\subsection{Por que gas existe}
O mecanismo de gas existe para:
\begin{itemize}[leftmargin=*]
  \item limitar execução infinita ou abusiva;
  \item cobrar pelo uso de armazenamento e computação compartilhada;
  \item incentivar contratos mais eficientes.
\end{itemize}

\subsection{Impacto no desenho de contratos}
Isso leva a consequências práticas:
\begin{itemize}[leftmargin=*]
  \item escrever em armazenamento custa mais do que ler;
  \item laços longos e estruturas desnecessárias aumentam custo;
  \item nem tudo deve ser colocado on-chain;
  \item cálculos pesados e grandes volumes de dados tendem a ser melhores fora da blockchain.
\end{itemize}

\section{Contratos inteligentes}
Um contrato inteligente é um programa armazenado na blockchain e executado segundo regras determinísticas da plataforma. Apesar do nome, ele não é ``inteligente'' no sentido humano; trata-se de código com regras explícitas.

Em uma aplicação didática de controle de acesso, um contrato pode ser responsável por:
\begin{itemize}[leftmargin=*]
  \item registrar o vínculo entre uma identidade e um identificador de referência;
  \item armazenar permissões ou janelas temporais de acesso;
  \item validar se um dispositivo está autorizado a consultar determinado estado;
  \item emitir eventos que possam ser auditados por outros componentes.
\end{itemize}

\section{Leitura, escrita e eventos}
Ao estudar contratos, vale separar três categorias:
\begin{itemize}[leftmargin=*]
  \item \textbf{funções de leitura}: consultam estado sem alterá-lo;
  \item \textbf{funções de escrita}: alteram estado e exigem transação;
  \item \textbf{eventos}: registros emitidos pelo contrato para facilitar observação externa.
\end{itemize}

\subsection{Funções de leitura}
São úteis para:
\begin{itemize}[leftmargin=*]
  \item consultar permissões atuais;
  \item recuperar apontadores ou identificadores vinculados;
  \item depurar o comportamento esperado do contrato.
\end{itemize}

\subsection{Funções de escrita}
São usadas quando a aplicação precisa:
\begin{itemize}[leftmargin=*]
  \item registrar novo estado;
  \item revogar ou conceder acesso;
  \item atualizar vínculos e políticas.
\end{itemize}

Como escrevem estado, essas funções têm custo, geram transações e exigem confirmação.

\subsection{Eventos}
Eventos não substituem armazenamento. Eles servem principalmente para:
\begin{itemize}[leftmargin=*]
  \item notificar sistemas externos;
  \item facilitar indexação e monitoramento;
  \item apoiar trilhas de auditoria.
\end{itemize}

Uma aplicação cliente ou um relayer pode escutar eventos e reagir a eles, mas a lógica principal de autorização não deve depender apenas de suposições implícitas sobre histórico de eventos se o estado atual pode ser consultado diretamente.

\section{Estruturas típicas em contratos}
Ao modelar contratos, é comum encontrar:
\begin{itemize}[leftmargin=*]
  \item variáveis de estado;
  \item \textit{mappings} entre identificadores e registros;
  \item estruturas compostas para armazenar atributos de autorização;
  \item verificações de permissões por remetente, papel ou contexto;
  \item eventos emitidos após mudanças relevantes.
\end{itemize}

Em termos de projeto, uma boa modelagem tenta equilibrar:
\begin{itemize}[leftmargin=*]
  \item simplicidade de verificação;
  \item custo on-chain;
  \item clareza de auditoria;
  \item separação entre dados mínimos on-chain e dados ricos off-chain.
\end{itemize}

\section{Armazenamento on-chain versus off-chain}
Nem todo dado deve ficar na blockchain. Há pelo menos três razões para isso:
\begin{itemize}[leftmargin=*]
  \item \textbf{custo}: armazenar muitos dados on-chain é caro;
  \item \textbf{privacidade}: alguns dados não devem ser públicos;
  \item \textbf{imutabilidade}: uma vez publicado, o dado on-chain é difícil de remover.
\end{itemize}

Por isso, arquiteturas modernas frequentemente colocam on-chain apenas:
\begin{itemize}[leftmargin=*]
  \item identificadores;
  \item hashes;
  \item apontadores;
  \item estados mínimos de autorização;
  \item evidências necessárias para verificação.
\end{itemize}

Os dados mais extensos ou sensíveis costumam ficar fora da blockchain, em serviços auxiliares ou armazenamentos distribuídos. Esse será o gancho do próximo capítulo sobre \textit{InterPlanetary File System} (IPFS).

\section{Por que blockchain faz sentido em controle de acesso}
Em um sistema didático de identidade e autorização, blockchain pode ser útil quando o objetivo é:
\begin{itemize}[leftmargin=*]
  \item registrar regras de acesso com trilha auditável;
  \item permitir verificação independente do estado autorizado;
  \item separar a política verificável da interface de administração;
  \item reduzir dependência de um único banco central opaco para decisão crítica.
\end{itemize}

Ela não elimina a necessidade de outros componentes. Firmware, relayer, frontend, armazenamento off-chain e gestão de chaves continuam sendo partes essenciais da arquitetura.

\section{Limitações e trade-offs}
Uma decisão madura sobre blockchain precisa reconhecer seus custos:
\begin{itemize}[leftmargin=*]
  \item latência maior do que sistemas puramente locais;
  \item custo financeiro ou computacional por transação;
  \item dificuldade de corrigir rapidamente dados mal modelados;
  \item exposição pública de parte do estado, conforme a rede usada;
  \item necessidade de desenhar com cuidado o que fica on-chain e o que fica fora.
\end{itemize}

Em outras palavras, blockchain não deve ser usada por moda. Ela deve ser escolhida quando os benefícios de verificabilidade, rastreabilidade e execução compartilhada superam seus custos operacionais.

\section{Relação com uma arquitetura como a do DIDGuard}
Em uma arquitetura inspirada no DIDGuard, o contrato inteligente tende a assumir o papel de camada verificável de política e vínculo mínimo. Isso significa que ele pode:
\begin{itemize}[leftmargin=*]
  \item associar um identificador descentralizado a um apontador externo;
  \item manter autorização por dispositivo, identidade ou janela temporal;
  \item fornecer funções de consulta de estado atual;
  \item emitir eventos úteis para observação do sistema.
\end{itemize}

Enquanto isso:
\begin{itemize}[leftmargin=*]
  \item o firmware lida com a credencial física e a coleta local;
  \item o relayer integra chamadas ao contrato e serviços auxiliares;
  \item o armazenamento distribuído hospeda dados mais ricos;
  \item o frontend fornece operação e visualização.
\end{itemize}

Essa divisão de responsabilidades é importante porque evita transformar o contrato em um repositório de tudo. O contrato deve guardar o que precisa ser público, verificável e mínimo.

\section{Leituras recomendadas}
\begin{itemize}[leftmargin=*]
  \item white paper do Bitcoin, para entender a noção original de livro-razão distribuído e encadeamento por hash \citep{nakamoto_bitcoin};
  \item white paper do Ethereum, para motivação da plataforma programável \citep{ethereum_whitepaper};
  \item \textit{Ethereum Yellow Paper}, para visão mais formal da execução e do custo computacional \citep{wood_ethereum};
  \item documentação oficial de Solidity, para modelagem de contratos, tipos, eventos e boas práticas \citep{solidity_docs};
  \item documentação de opcodes da EVM, para aprofundamento no modelo de execução \citep{evm_codes}.
\end{itemize}

\section{Rodando na prática}
O objetivo desta prática é implantar um contrato simples em uma blockchain local, gravar um estado de autorização e observar a diferença entre leitura, escrita e evento emitido.

\subsection{Pré-requisitos}
\begin{itemize}[leftmargin=*]
  \item Node.js 18 ou superior;
  \item terminal com suporte a \texttt{npm} e \texttt{npx};
  \item diretório \texttt{docs/book/examples/cap4-blockchain-contratos} disponível localmente.
\end{itemize}

\subsection{Arquivos do laboratório}
O mini projeto deste capítulo contém:
\begin{verbatim}
contracts/AccessRegistry.sol
scripts/deploy.js
scripts/demo.js
hardhat.config.js
package.json
\end{verbatim}

\subsection{Instalação}
No diretório do laboratório:
\begin{verbatim}
npm install
\end{verbatim}

\subsection{Execução}
Em um terminal:
\begin{verbatim}
npx hardhat node
\end{verbatim}

Em outro terminal:
\begin{verbatim}
npx hardhat run scripts/deploy.js --network localhost
\end{verbatim}

Copie o endereço do contrato implantado e execute:
\begin{verbatim}
$env:ACCESS_REGISTRY_ADDRESS="COLE_O_ENDERECO_AQUI"
npx hardhat run scripts/demo.js --network localhost
\end{verbatim}

Alternativamente, para executar deploy e demonstração em um único passo (sem copiar endereços):
\begin{verbatim}
npx hardhat run scripts/demo-auto.js --network localhost
\end{verbatim}

\subsection{O que observar}
\begin{itemize}[leftmargin=*]
  \item o contrato sendo implantado na rede local;
  \item a chamada de escrita \texttt{setAccess} gerando uma transação;
  \item a chamada de leitura \texttt{getAccess} retornando o estado sem criar nova transação;
  \item o evento \texttt{AccessUpdated} sendo recuperado do recibo.
\end{itemize}

\subsection{Perguntas de análise}
\begin{itemize}[leftmargin=*]
  \item Por que a leitura pode ser modelada como \texttt{view}?
  \item Por que a escrita precisa de uma transação assinada e minerada?
  \item Qual informação foi mantida on-chain e qual informação ficou fora do contrato?
\end{itemize}

\subsection{Extensão sugerida}
Evolua o contrato para:
\begin{itemize}[leftmargin=*]
  \item registrar o identificador de um dispositivo;
  \item armazenar um tempo de expiração;
  \item impedir atualizações com prazo já expirado.
\end{itemize}

\subsection{Localização do exemplo}
Todos os arquivos usados nesta prática estão em:
\begin{verbatim}
docs/book/examples/cap4-blockchain-contratos
\end{verbatim}

\vspace{1em}
\noindent\rule{\textwidth}{0.4pt}
\vspace{0.5em}
\noindent\textbf{Encerramento e próximos passos.}
Com este capítulo, a apostila cobriu as quatro camadas fundamentais do DIDGuard: identidade descentralizada (DIDs), credencial física (RFID/MIFARE), primitivas criptográficas e registro verificável de permissões (blockchain). Os capítulos seguintes --- quando disponíveis --- abordarão armazenamento distribuído com IPFS, firmware do ESP32, integração via relayer e frontend de operação, completando o percurso da teoria à implementação.
